% Modelo base — Currículo Camaleão
% Use o preâmbulo abaixo SEM alterar a estrutura. Substitua apenas o conteúdo
% entre os marcadores de seção (%% =========================).
% Regras de formatação por trás deste preâmbulo: ver regras-ats.md.

\documentclass[10pt,a4paper]{article}

\usepackage[utf8]{inputenc}
\usepackage[T1]{fontenc}
\usepackage[brazil]{babel}
\usepackage{lmodern}
\usepackage[margin=1.8cm]{geometry}
\usepackage{enumitem}
\usepackage{titlesec}
\usepackage{hyperref}

\hypersetup{
    colorlinks=true,
    urlcolor=black
}

\pagestyle{empty}
\setlength{\parindent}{0pt}
\setlength{\parskip}{0.4em}
\setlist[itemize]{leftmargin=1.3em, itemsep=0.2em}

\titleformat{\section}
  {\large\bfseries}
  {}
  {0pt}
  {}
  []

\titlespacing*{\section}{0pt}{1em}{0.3em}

\begin{document}

\textbf{\Large Cauê Franco} \\
AI Full Stack Engineer (Agentes de LLM | Python | TypeScript)

João Pessoa - Paraíba, Brasil \\
Telefone: +55 83 98746-8188 \\
Email: cauefranco01@gmail.com \\
LinkedIn: https://www.linkedin.com/in/cauefranco01 \\
GitHub: https://github.com/cauedev30

% =========================
\section*{Resumo Profissional}

AI Full Stack Engineer, construo agentes de LLM e o produto inteiro em volta
deles: back-end, interface, integrações e deploy. Tenho experiência com
agentes de IA e sistemas multiagentes, Retrieval-Augmented Generation (RAG),
grounding em base de conhecimento, function calling, MCP servers e evals com
controle de custo de tokens. Desenvolvo em Python (FastAPI, SQLAlchemy) e
TypeScript (React, Next.js, Node.js) sistemas que rodam em produção: construí
de ponta a ponta um agente de atendimento no WhatsApp que tornou 100\%
automática a triagem antes manual, e automatizei o provisionamento de cerca de
20 contas Chatwoot com n8n e Supabase. Trabalho com workflow agêntico (Claude
Code com skills e hooks), PostgreSQL, Docker e CI/CD com GitHub Actions.

% =========================
\section*{Experiência Profissional}

\textbf{GeniAI} \\
Cargo: Generative AI Engineer (efetivado em 07/2026; antes, Especialista em Automação, estágio) \\
Período: 11/2025 -- Presente \\
Local: João Pessoa, PB

\begin{itemize}
\item Construí de ponta a ponta um agente LLM de atendimento no WhatsApp que
identifica o cliente, abre o ticket e responde com grounding em uma base de
conhecimento curada, entendendo áudio e imagem, e o produto em volta dele:
back-end em Python, FastAPI e PostgreSQL, front-end em React e Next.js com
kanban de tickets e painel de monitoramento e indicadores, integração com a
API da OpenAI e testes automatizados com Pytest e Vitest; a triagem, que era
manual, passou a ser 100\% automática.
\item Automatizei o provisionamento de atendentes em cerca de 20 contas
Chatwoot de uma rede de franquias de saúde com n8n e Supabase (PostgreSQL):
fluxo idempotente e parametrizado, via APIs REST e webhooks, que cria usuário,
times, etiquetas, regras de roteamento e inboxes, confere banco e Chatwoot
antes de escrever e registra cada execução.
\item Aplico workflow agêntico como método: harness em volta do modelo
(skills, hooks, guardrails, evals e contexto versionado), com cada decisão
técnica ligada ao planejamento do projeto e a um resultado de negócio.
\end{itemize}

% STACK DE RESERVA: agente de atendimento (fato confirmado pelo Cauê, 22/10/2026).
% Não aparece no PDF. Use só o item que a vaga pedir, no bullet do agente e em
% Competências; sem pedido da vaga, fica de fora.
%   Back-end: Python 3.12, FastAPI, Uvicorn
%   Banco: PostgreSQL com SQLAlchemy async + asyncpg
%   Front-end (quadro de tickets): Next.js 16, React 19, TypeScript
%   Arrastar cards no kanban: dnd-kit
%   LLM: OpenAI (modelo gpt-6-luna) chamada via httpx
%   Áudio: transcrição (speech-to-text) pela OpenAI
%   Atendimento: Chatwoot + WAHA (API de WhatsApp)
%   Testes: Pytest, Vitest, Testing Library, Playwright
%   Qualidade: ruff, mypy strict, ESLint
%   Deploy: Docker Compose em servidor próprio + túnel Cloudflare (Cloudflare Tunnel)

\textbf{Desenvolvedor Autônomo / Freelancer} \\
Cargo: Desenvolvedor Full Stack \\
Período: 07/2026 -- Presente \\
Local: Remoto, João Pessoa, PB

\begin{itemize}
\item Desenvolvi um site de pedidos para restaurante em Next.js (JavaScript),
publicado no edge da Cloudflare com Cloudflare Workers.
\item Implementei painel autenticado para o cliente organizar os pedidos,
otimizando o fluxo de trabalho da operação.
\end{itemize}

% STACK DE RESERVA: site de pedidos do restaurante (fato confirmado pelo Cauê,
% 22/10/2026). Não aparece no PDF; o nome do cliente nunca entra no currículo.
% Use só o item que a vaga pedir, no bullet do freelance e em Competências.
%   Framework: Next.js 16 (App Router) + React 19
%   Linguagem: JavaScript
%   Estilo: CSS puro com design tokens e CSS Modules
%   Fontes e imagens: next/font e next/image
%   Estado: React Context (carrinho de compras)
%   Back-end: API Routes do Next.js (ex.: rota do cardápio do dia)
%   Checkout: pedido montado no site e enviado pelo WhatsApp
%   Testes: Vitest + Testing Library (397 testes automatizados)
%   Qualidade: ESLint
%   Deploy: Cloudflare Workers via OpenNext + Wrangler

\textbf{Fábrica de Software UBTech Office (Unipê)} \\
Cargo: Desenvolvedor Back-end (estágio) \\
Período: 08/2025 -- 11/2025 \\
Local: João Pessoa, PB

\begin{itemize}
\item Trabalhei no desenvolvimento do back-end de um sistema de gestão
operacional para um órgão governamental da Paraíba, utilizando Python e Django
Framework, centralizando informações críticas.
\item Modelei e estruturei bancos de dados relacionais com PostgreSQL e
SQLite3, garantindo integridade, segurança e performance no acesso aos dados
institucionais.
\end{itemize}

% =========================
\section*{Competências Técnicas}

\textbf{IA e Agentes:} Agentes de IA e sistemas multiagentes, LLMs,
Retrieval-Augmented Generation (RAG), embeddings, bancos de dados vetoriais,
grounding, function calling e tool use, MCP (Model Context Protocol), prompt e
context engineering, evals, controle de custo de tokens, entrada multimodal
(áudio e imagem), API da OpenAI \\
\textbf{Automação e Integração:} n8n, Make, Chatwoot, APIs REST, webhooks \\
\textbf{Linguagens:} Python, TypeScript, JavaScript \\
\textbf{Full Stack:} FastAPI, SQLAlchemy, Django, Flask, Node.js, React,
Next.js, Clean Architecture \\
\textbf{Banco de Dados:} PostgreSQL, SQLite, Supabase (Row Level Security) \\
\textbf{Cloud e DevOps:} Docker, Dokku, Cloudflare Workers, GitHub Actions
(CI/CD), Git, GitHub \\
\textbf{Testes:} Pytest, Vitest, Playwright \\
\textbf{Ferramentas e Metodologias:} critério de aceite executável,
Claude Code: workflow agêntico (harness, guardrails, contexto versionado,
evals, policies)

% =========================
\section*{Formação Acadêmica}

\textbf{Centro Universitário Unipê} \\
Curso: Análise e Desenvolvimento de Sistemas \\
Período: 2025 -- 2027 \\
Status: Em andamento

% =========================
\section*{Idiomas}

Português: Nativo \\
Inglês: Intermediário/Avançado

\end{document}
